\section{Introduction}
\begin{definition}
    A topology on a set $X$ is a collection $\tau$ of subsets of $X$ such that:
    \begin{enumerate}\label{1.1}
        \item $\emptyset,X \in \tau$.
        \item If $U_i \in \tau$ for all $i \in I$, then $\bigcup_{i \in I} U_i \in \tau$.
        \item If $U_i \in \tau$ for all $i \in I$ where $I$ is finite, then $\bigcap_{i \in I} U_i \in \tau$.
    \end{enumerate}
    We call the pair $(X,\tau)$ a topological space. The elements of $\tau$ are called open sets.
\end{definition}
\begin{definition}\label{1.2}
    A \textbf{basis} for a topology on $X$ is a collection $\mathcal{B}$ of subsets of $X$ such that:
    \begin{enumerate}
        \item $\forall x \in X$, $\exists B\in \mathcal{B}$ containing $x$.
        \item If $x \in B_1 \cap B_2$ for some $B_1,B_2 \in \mathcal{B}$, then there is a basis element $B_3 \in \mathcal{B}$ such that $x \in B_3 \subseteq B_1 \cap B_2$.
    \end{enumerate}
    $\mathcal{B}$ generates a topology $\tau$ on $X$ as follows: \\
    $U \in \tau$ if $\forall x \in U$, $\exists B \in \mathcal{B}$ such that $x \in B \subseteq U$.
\end{definition}
\begin{remark}\label{1.3}
    It is important that $B\in \mathcal{B}$ is a subset of $U$ in the last statement. If we only required $x \in B \cap U$, we would run into problems. For example consider $X = \mathbb{R}$ and $U=\{0,1\}$. We know that $\mathcal{B}=\{(a,b):a,b\in \mathbb{R}, a<b\}$ is a basis for the standard topology on $\mathbb{R}$. Then for $x=0$, we can choose $B=(-1,1)$ and for $x=1$, we can choose $B=(0,2)$. This would mean that $U$ is open, which should not be the case. 
\end{remark}
\begin{example}\label{1.4}
    \begin{itemize}
        \item $(M,d)$ is a metric space. Then $\mathcal{B}=\{B_r(x):x\in M, r>0\}$ is a basis for the metric topology on $M$, where $B_r(x)=\{y\in M:d(x,y)<r\}$ is the open ball of radius $r$ around $x$.
        \item $X,Y$ topological spaces. The product topology on $X\times Y$ is generated by the basis $\mathcal{B}=\{U\times V:U\subseteq X, V\subseteq Y \text{ open}\}$.
        \item subspace topology: Let $(X,\tau)$ be a topological space, $S\subseteq X$. Then the subspace topology of $S$ is $\tau_S=\{U\cap S:U\in \tau\}$.
    \end{itemize}
\end{example}
\begin{definition}\label{1.5}
    $U$ is a neighbourhood of $x$ if $\exists V$ open such that $x \in V \subseteq U$.
\end{definition}
\begin{definition}\label{1.6}
    Let $X,Y$ be topological spaces. \\
    A map $f:X\to Y$ is continuous if $\forall U\subseteq Y$ open, $f^{-1}(U)$ is open in $X$.\\
    If $f$ is a bijection and both $f$ and $f^{-1}$ are continuous, then $f$ is called a homeomorphism. In this case, we say that $X$ and $Y$ are homeomorphic or topologically equivalent ($X\simeq Y$).
\end{definition}
\begin{example}\label{1.7}
    \begin{itemize}
        \item  $f:[0,1]\to \mathbb{S}^1$ given by $f(t) = (cos(2\pi t), sin(2\pi t))$. \\
        $f$ is continuous and bijective but $f^{-1}$ is not continuous.\\
        $f([0,1))$ is not open in $\mathbb{S}^1$.
        \item $\mathbb{R}\not\simeq \mathbb{R}^2$: Assume they are homeomorphic. Then there exists a homeomorphism $f:\mathbb{R}\to \mathbb{R}^2$ with $f(0)=(a,b)$. But $f|_{\mathbb{R}\setminus\{0\}}:\mathbb{R}\setminus\{0\}\to \mathbb{R}^2\setminus\{(a,b)\}$ is a homeomorphism. But $\mathbb{R}\setminus\{0\}$ is not connected, while $\mathbb{R}^2\setminus\{(a,b)\}$ is connected. This is a contradiction. 
        \item $\mathbb{R}^2\not\simeq \mathbb{R}^3$: Assume they are diffeomorphic. Then there exists a diffeomorphism $f:\mathbb{R}^2\to \mathbb{R}^3$. Then $f\circ f^{-1}=id_{\mathbb{R}^3}, f^{-1}\circ f=id_{\mathbb{R}^2}$ so $Df_xDf^{-1}_y=I_3$ and $Df^{-1}_yDf_x=I_2$. 
    \end{itemize}
\end{example}